\leftskip 0in
\parindent 1.27cm % Sangría según el manual (1,27 cm): la volvemos a establecer.
\newpage
\section{Anexos}

% =====================================================================
\subsection{Reflexión personal}
% =====================================================================
% Texto propio del practicante. Sobre la versión anterior solo se corrigieron
% tildes, concordancias y una repetición literal de frase; el estilo se mantuvo
% tal cual estaba escrito.

Cuando llegué a DTS, lo primero que recibí fue una clase por parte de mi supervisor. En un par de horas recibí detalles técnicos de cómo tenía que trabajar durante los próximos dos meses. Sin embargo, esta dinámica de no saber nada, investigar sobre el tema, para luego crear diagramas, tablas y documentación fue lo que definió el resto de mi práctica.

Como practicante, tuve que crear un proyecto prácticamente desde cero, buscar documentación, aprender sobre el flujo de trabajo, qué se espera de los documentos que estaba creando, etc. Gran parte del aprendizaje fue durante la marcha y tuve que adaptarme a este entorno de trabajo, donde cada pocas semanas tenía que mostrar todo lo que aprendí y apliqué al proyecto.

En particular, el hecho de tener reuniones periódicas con mi jefe y la contraparte técnica de INVAP me exigía mantener un ritmo constante y actualizar la documentación cada pocos días. Esto me obligó a buscar una metodología de trabajo que me permitiera trazar con facilidad las decisiones de diseño. A medida que el proyecto fue creciendo conmigo, llegó un punto en donde se volvió muy difícil mantener la trazabilidad de los requerimientos y las decisiones de diseño, por lo que propuse migrar el proyecto a ARCADIA/Capella. Esto me permitió mantener la trazabilidad de manera más eficiente y me enseñó a tomar decisiones técnicas que impactan en el proyecto.

Por otra parte, el hecho de tener que buscar mis propias referencias para el desarrollo del proyecto facilitó mi flujo de trabajo, ya que me ayudó a aclarar cuáles eran los puntos importantes en cada entrega. Gracias a esto pude estandarizar la documentación y mantener un ritmo de trabajo constante, lo que finalmente me permitió cumplir con los objetivos de la práctica.

Ahora bien, la universidad me entregó una base teórica que le dio validez a los resultados obtenidos en mi práctica. Por ejemplo, los cursos de comunicaciones, como Comunicaciones Inalámbricas, me permitieron entender cómo funciona el enlace de comunicación y cómo se comporta la señal en el espacio. Por otra parte, los cursos de programación fueron los que me dieron las herramientas para crear un simulador de enlace con el que pude validar los resultados obtenidos en la práctica. Además, cursos como Capstone me ayudaron a entender cómo se desarrolla un proyecto de ingeniería y cómo se debe documentar.

Sin embargo, todo lo relacionado a la gestión de proyectos de ingeniería, la ingeniería de sistemas y la gestión de riesgos no fue parte de mi formación en la universidad. Esto me obligó a aprender sobre la marcha, buscando referencias y aplicando lo aprendido en el proyecto, y fue probablemente el aprendizaje más valioso que me llevo de estos dos meses.

En general, la práctica me permitió aplicar los conocimientos adquiridos en la universidad y aprender sobre la marcha, lo que me ayudó a desarrollar habilidades técnicas y de gestión de proyectos. También me permitió entender qué habilidades son necesarias para desarrollarse en el mundo laboral y cómo se deben aplicar los conocimientos adquiridos en la universidad. Aspectos como la creatividad, la proactividad, la capacidad de aprender y adaptarse a nuevas situaciones, y la capacidad de trabajar en equipo son habilidades que se deben desarrollar para poder desempeñarse en el mundo laboral.

\newpage
% ARCHIVO GENERADO AUTOMÁTICAMENTE por content/gen_tablas.py
% No editar a mano: vuelve a correr el script si cambia el Excel.

\subsection{Anexo A: requerimientos de usuario}

\begin{footnotesize}
\begin{longtable}{|@{\hspace{4pt}}l >{\raggedright\arraybackslash}p{5.4cm} c c c >{\raggedright\arraybackslash}p{4.6cm}@{\hspace{4pt}}|}
\caption{Requerimientos de usuario del proyecto BIFROST. Elaboración propia.\vspace{1ex}}\label{tab:anexoA}\\
\toprule
\textbf{ID} & \textbf{Requerimiento} & \textbf{Prior.} & \textbf{Riesgo} & \textbf{Tipo} & \textbf{Mitigación propuesta} \\
\midrule\endfirsthead
\multicolumn{6}{@{}l@{}}{\textit{(continuación)}}\\
\toprule
\textbf{ID} & \textbf{Requerimiento} & \textbf{Prior.} & \textbf{Riesgo} & \textbf{Tipo} & \textbf{Mitigación propuesta} \\
\midrule\endhead
\bottomrule\endfoot
UR-01 & El sistema deberá proveer un servicio \textit{Store \& Forward} satelital, almacenando datos durante su órbita para ser descargados en una estación terrestre. & Alta & Alto & Funcional & Diseño de memoria no volátil con margen $\geq$50\%, simulación de carga orbital y pruebas de estrés en tierra \\
UR-02 & El sistema deberá soportar un enlace \textit{downlink} desde el satélite hacia estaciones terrenas en territorio chileno. & Alta & Medio & Funcional &  \\
UR-03 & El sistema deberá soportar un enlace \textit{uplink} desde estaciones remotas de monitoreo meteorológico, sismológico y/o vulcanológico hacia el satélite para la recolección de datos. & Alta & Medio & Funcional & Validación temprana con transmisores representativos y pruebas HIL \\
UR-04 & El sistema deberá ser compatible con el estándar internacional para recolección de datos ambientales por satélite utilizado por estaciones remotas autónomas de meteorología, sismología, vulcanología, etc. & Alta & Medio & Proceso & Revisión de estándar ARGOS y asesoría INVAP \\
UR-05 & El sistema deberá asegurar la integridad de los datos durante su recepción, almacenamiento y transmisión. & Media & Bajo & Funcional &  \\
UR-06 & El sistema deberá ser compatible con plataforma CubeSat 12U en órbita terrestre baja (LEO) sol síncrona a una altitud nominal aproximada de 650 km con una inclinación de 96$^{\circ}$. & Alta & Alto & No funcional & Análisis de interfaces mecánicas/eléctricas y revisión PDR con proveedor de plataforma \\
UR-07 & El sistema deberá cubrir el territorio chileno continental, zonas extremas, territorio insular y territorio antártico. & Media & Bajo & Funcional &  \\
UR-08 & El sistema deberá permitir el uso de componentes COTS siempre que cumplan con: (a) funcionamiento en órbita probado, (b) tolerancia a la radiación verificada, o (c) mitigaciones de diseño que reduzcan el riesgo a un nivel aceptable para la misión. & Alta & Alto & No funcional & Selección COTS con validación espacial y mitigación por redundancia y watchdogs \\
UR-09 & El sistema deberá ser de carácter experimental, con proyección a futuras aplicaciones comerciales. & Baja & Bajo & No funcional &  \\
UR-10 & El sistema deberá presentar una arquitectura modular, reutilizable y escalable. & Baja & Bajo & No funcional &  \\
UR-11 & El sistema debe garantizar una vida útil operativa mínima de 5 años en el entorno de radiación LEO definido. & Alta & Alto & No funcional & Pruebas de estrés, análisis térmico y pruebas de vida aceleradas \\
UR-12 & El sistema debe ser validado mediante pruebas funcionales en tierra y demostración en órbita. & Baja & Bajo & Proceso &  \\
UR-13 & El sistema deberá ser compatible con el sistema de separación del lanzador seleccionado EXOpod 12U, cumpliendo con las restricciones geométricas, mecánicas y de seguridad. &  &  & No funcional &  \\
\end{longtable}
\vspace{-1ex}\noindent\parbox{\textwidth}{\footnotesize Nota. UR-13 se incorporó al cierre de la PDR, con posterioridad al análisis de riesgos, por lo que todavía no tiene prioridad ni nivel de riesgo asignados.}
\end{footnotesize}

\subsection{Anexo B: requerimientos de sistema}

\begin{footnotesize}
\begin{longtable}{|@{\hspace{4pt}}l >{\raggedright\arraybackslash}p{5.4cm} c c c >{\raggedright\arraybackslash}p{4.6cm}@{\hspace{4pt}}|}
\caption{Requerimientos de sistema del proyecto BIFROST. Elaboración propia.\vspace{1ex}}\label{tab:anexoB}\\
\toprule
\textbf{ID} & \textbf{Requerimiento} & \textbf{Prior.} & \textbf{Riesgo} & \textbf{Tipo} & \textbf{Mitigación propuesta} \\
\midrule\endfirsthead
\multicolumn{6}{@{}l@{}}{\textit{(continuación)}}\\
\toprule
\textbf{ID} & \textbf{Requerimiento} & \textbf{Prior.} & \textbf{Riesgo} & \textbf{Tipo} & \textbf{Mitigación propuesta} \\
\midrule\endhead
\bottomrule\endfoot
SR-01 & El sistema deberá permitir la recepción de paquetes de datos con un tamaño desde 32 a 256 bits, con una tasa de datos de 400 bps. & Alta & Alto & Funcional & Prototipo temprano del receptor y pruebas BER en laboratorio \\
SR-02 & El sistema deberá operar como receptor en el rango de frecuencias 401 - 403 MHz utilizado por estaciones remotas, con una polarización circular derecha y un ancho de haz de 130$^{\circ}$ HPBW. & Alta & Alto & No funcional & Verificación de front-end RF y filtros certificados \\
SR-03 & El sistema deberá ser compatible con la estructura de datos enviado por las estaciones remotas, incluyendo modulación, canalización y protocolos, además de las variaciones en frecuencia por efecto Doppler. & Alta & Alto & Funcional & Simulación Doppler y pruebas con generador de señales \\
SR-04 & El sistema deberá almacenar al menos 1.000.000 de mensajes durante un tiempo máximo de 15 órbitas y sin pérdidas hasta establecer contacto con una estación terrestre para la descarga. & Media & Medio & Funcional & Margen de memoria $\geq$50\% y control de overflow \\
SR-05 & El sistema deberá mantener la asociación temporal de los datos almacenados. & Baja & Bajo & Funcional &  \\
SR-06 & El sistema deberá permitir la descarga de al menos 10.000 mensajes almacenados hacia estaciones terrestres cuando exista conectividad disponible. & Media & Medio & Funcional & Ensayos de descarga en peor escenario \\
SR-07 & El sistema deberá entregar los datos al subsistema de comunicaciones de la plataforma (\textit{Transponder} Banda S) para el enlace \textit{downlink}. & Media & Medio & Proceso &  \\
SR-08 & El sistema deberá comunicarse con el OBC del CubeSat 12U mediante el bus definido por la plataforma (CANBUS), gestionando una tasa de transferencia de datos compatible y reportando periódicamente su estado operativo. & Alta & Alto & No funcional & Pruebas de compatibilidad CAN con OBC real \\
SR-09 & El sistema deberá operar bajo el control del OBC, escuchando activamente en las zonas definidas por la misión y transmitiendo únicamente cuando sea ordenado. & Media & Bajo & Funcional & Validación de lógica de estados y temporización \\
SR-10 & El sistema deberá implementar mecanismos de integridad de datos en todas las etapas, con un error de bit menor a 0,01\% y control de acceso a la descarga de datos por estaciones terrestres. & Media & Medio & Interfaz & CRC, ACK/NACK y control de acceso \\
SR-13 & El sistema deberá consumir un máximo de 10 W en operación nominal y ser alimentado por el sistema de potencia del satélite de forma continua. & Alta & Alto & No funcional & Presupuesto de potencia con margen $\geq$25\% \\
SR-14 & El sistema deberá tener dimensiones máximas 20 cm x 20 cm x 10 cm. & Alta & Alto & No funcional & Revisión de \textit{layout} 3D y \textit{mockup} mecánico \\
SR-15 & El sistema deberá tener una masa máxima de 3,0 kg y mantener su Centro de Gravedad (CoG) a $\pm$2 mm del centro geométrico, el cual coincide con el eje longitudinal geométrico del satélite. & Alta & Alto & No funcional & Simulaciones CAD para estimaciones de masas y revisión periódica de masa en el prototipo \\
SR-16 & El sistema deberá cumplir con requisitos de EMI/EMC compatibles con operación satelital, evitando interferencias internas y externas durante todo el proceso de lanzamiento, despliegue y operación. & Alta & Alto & No funcional & Mediciones EMI/EMC y separación RF \\
SR-17 & El sistema deberá diseñarse de forma modular y permitir la actualización de \textit{firmware} en vuelo (OTA). & Media & Bajo & No funcional &  \\
SR-18 & El sistema deberá certificarse mediante pruebas termo-vacío, mecánicas, EMI y eléctricas en las instalaciones de Centro Espacial Nacional y ser validado en órbita. & Media & Bajo & Proceso &  \\
SR-19 & El sistema deberá respetar la envolvente interna permitida, los volúmenes prohibidos y las restricciones sobre protrusiones, mecanismos móviles y bordes activos, asegurando la integración con la plataforma. &  &  & No funcional &  \\
SR-20 & El sistema deberá soportar los ambientes mecánicos de lanzamiento (cargas cuasiestáticas, vibración y shock) definidos por el lanzador y \textit{deployer} seleccionados. &  &  & No funcional &  \\
SR-21 & El sistema deberá implementar \textit{Watchdog Timers} externos e independientes para reinicio ante eventos SEU/\textit{Latch-up}. &  &  & Funcional &  \\
SR-22 & Protección electrónica contra \textit{Single Event Latch-up} (SEL) mediante limitación de corriente activa. &  &  & No funcional &  \\
SR-23 & Los componentes y materiales no metálicos que utilice el sistema deberán cumplir con un TML $<$ 1,0\% y un CVCM $<$0,1\% &  &  & No funcional &  \\
\end{longtable}
\vspace{-1ex}\noindent\parbox{\textwidth}{\footnotesize Nota. La numeración salta de SR-10 a SR-13 porque SR-11 y SR-12 se eliminaron al consolidar la línea base y sus identificadores no se reutilizaron. SR-19 a SR-23 se incorporaron al cierre de la PDR y todavía no tienen evaluación de riesgo. SR-07 conserva la redacción aprobada en la SRR, cuando la línea base contemplaba el uso del transponder en banda S de la plataforma; durante la PDR el enlace de bajada se reasignó a una antena UHF propia de la carga útil, que es la configuración que describe el cuerpo de este informe. La actualización formal del requerimiento queda pendiente para la CDR.}
\end{footnotesize}

\subsection{Anexo C: matriz de trazabilidad de requerimientos}

\begin{footnotesize}
\begin{longtable}{|@{\hspace{4pt}}l >{\raggedright\arraybackslash}p{12.2cm}@{\hspace{4pt}}|}
\caption{Matriz de trazabilidad entre requerimientos de usuario y de sistema. Elaboración propia.\vspace{1ex}}\label{tab:anexoC}\\
\toprule
\textbf{Req. de usuario} & \textbf{Requerimientos de sistema asociados} \\
\midrule\endfirsthead
\multicolumn{2}{@{}l@{}}{\textit{(continuación)}}\\
\toprule
\textbf{Req. de usuario} & \textbf{Requerimientos de sistema asociados} \\
\midrule\endhead
\bottomrule\endfoot
UR-01 & SR-04, SR-06 \\
UR-02 & SR-03, SR-07 \\
UR-03 & SR-02, SR-03 \\
UR-04 & SR-01, SR-02, SR-03 \\
UR-05 & SR-04, SR-05, SR-08, SR-10 \\
UR-06 & SR-13, SR-14, SR-15, SR-16, SR-19, SR-20 \\
UR-07 & SR-09, SR-10 \\
UR-08 & SR-16, SR-17, SR-18, SR-21, SR-23 \\
UR-09 & SR-17 \\
UR-10 & SR-17 \\
UR-11 & SR-17, SR-18, SR-21, SR-23 \\
UR-12 & SR-18 \\
UR-13 & SR-19, SR-22 \\
\end{longtable}
\vspace{-1ex}\noindent\parbox{\textwidth}{\footnotesize Nota. La matriz registra 35 relaciones. Los 21 requerimientos de sistema aparecen al menos una vez, de modo que la cobertura es completa en ambos sentidos.}
\end{footnotesize}

\subsection{Anexo D: matriz preliminar de ensayos de verificación}

\begin{footnotesize}
\begin{longtable}{|@{\hspace{4pt}}l >{\raggedright\arraybackslash}p{4.0cm} >{\raggedright\arraybackslash}p{8.6cm}@{\hspace{4pt}}|}
\caption{Matriz preliminar de ensayos de verificación. Elaboración propia.\vspace{1ex}}\label{tab:anexoD}\\
\toprule
\textbf{ID} & \textbf{Nombre del ensayo} & \textbf{Descripción} \\
\midrule\endfirsthead
\multicolumn{3}{@{}l@{}}{\textit{(continuación)}}\\
\toprule
\textbf{ID} & \textbf{Nombre del ensayo} & \textbf{Descripción} \\
\midrule\endhead
\bottomrule\endfoot
FUN-01 & Caracterización de Potencia & Medición de consumo de corriente en todos los modos (Sleep, RX, Processing, TX Interface). Verificación de picos de arranque (\textit{Inrush Current}). \\
FUN-02 & Sensibilidad y Selectividad RF & Verificación del umbral de recepción (-115 dBm estimado) y ancho de banda en 401-403 MHz usando generador de señales. \\
FUN-03 & Simulación de Efecto Doppler & Inyección de señales con desplazamiento de frecuencia variable ($\pm$10 kHz) para validar el algoritmo de corrección del receptor. \\
FUN-04 & Integridad de Datos (BER) & Transmisión de 1 millón de paquetes y verificación de tasa de error de bit (BER) y corrección CRC tras almacenamiento. \\
FUN-05 & Interfaz CAN Bus \& CSP & Validación de intercambio de comandos y telemetría con emulador de OBC (o OBC real). Verificación de manejo de errores de bus. \\
FUN-06 & Prueba de \textit{Watchdog} y \textit{Failsafe} & Inyección de fallos de \textit{software} y \textit{hardware} (reset forzado) para medir tiempo de recuperación y estado seguro. \\
FUN-07 & Actualización de Firmware OTA & Simulación de carga de nueva imagen de \textit{software} a través del bus de datos y verificación de \textit{checksum} e instalación. \\
THR-01 & Ciclado Térmico (Presión Ambiente) & Ciclos de temperatura entre -30$^{\circ}$C y +70$^{\circ}$C en cámara climática para detectar fallas de soldadura o componentes infantiles (\textit{Stress screening}). \\
THR-02 & Vacío Térmico (TVAC) - \textit{Bake-out} & Exposición a alto vacío ($<$10$^{-5}$ Torr) y temperatura elevada (+60$^{\circ}$C) para desgasificación (\textit{Outgassing}) y limpieza. \\
THR-03 & Balance Térmico en Vacío & Validación del modelo térmico: verificación de temperaturas de operación en vacío (solo conducción) en casos fríos y calientes. \\
MEC-01 & Propiedades de Masa & Medición precisa de masa total ($<$ 3kg) y determinación del Centro de Gravedad (CoG) y Momentos de Inercia (MoI). \\
MEC-02 & Vibración Aleatoria (\textit{Random Vib}) & Ensayo en \textit{shaker} electrodinámico según perfiles de lanzamiento (ej. NASA GEVS o perfil Falcon 9). Verificación de integridad estructural. \\
MEC-03 & Vibración Sinusoidal (\textit{Sine Sweep}) & Barrido de frecuencias para identificar resonancias naturales (\textit{Eigenfrequencies}) y asegurar que $>$ 100 Hz (típico CubeSat). \\
MEC-04 & Ensayo de Choque (\textit{Shock}) & Simulación de cargas de choque por separación (si aplica por requerimiento de lanzador) o despliegue de antenas. \\
MEC-05 & Despliegue de Antenas & Verificación funcional del mecanismo de despliegue de antenas UHF (\textit{Burn-wire} / Resorte) en condiciones de gravedad compensada. \\
EMC-01 & Emisiones Conducidas (CE) & Medición de ruido inyectado por la carga útil hacia las líneas de potencia del bus (\textit{Power Bus Noise}). \\
EMC-02 & Susceptibilidad Radiada (RS) & Verificación de operación nominal ante campos externos (ej. Radares de lanzamiento o transmisor Banda S propio). \\
EMC-03 & Emisiones Radiadas (RE) & Medición de ruido RF generado por la electrónica digital (relojes, CPU) para asegurar ''silencio'' en banda GPS y Banda S receptor. \\
EMC-04 & Auto-Compatibilidad & Verificación de que la transmisión UHF propia no resetea la CPU ni afecta sensores cercanos. \\
\end{longtable}
\vspace{-1ex}\noindent\parbox{\textwidth}{\footnotesize Nota. Los 19 ensayos corresponden a especificaciones de verificación definidas durante la PDR y ninguno se ejecutó dentro del período de la práctica. Las menciones a banda S se refieren a los equipos de la plataforma, que actúan como fuente de interferencia sobre la carga útil.}
\end{footnotesize}

